% ============================================================
%  Multiplicação e Divisão de Inteiros — Apostila Premium | PratiqueJá
%  Engine: XeLaTeX
% ============================================================
\documentclass[12pt]{article}
\usepackage{../../pratiqueja}

% \hlm — destaque de resultado em modo-math (cor sem bold)
\newcommand{\hlm}[1]{{\color{darkblue}#1}}

\setsubject{Multiplicação e Divisão de Inteiros}

\begin{document}
\pagenumbering{arabic}
\setstretch{1.2}
\color{bodytext}

\heroheader{Multiplicação e Divisão de Inteiros}%
           {Regra dos Sinais, Potências e Aplicações}%
           {Matemática · Básico}

\setlength{\columnsep}{18pt}
\setlength{\columnseprule}{0.5pt}
\def\columnseprulecolor{\color{exborder}}

\begin{multicols}{2}

%─────────────────────────────────────────────────────────────
\TW{?}{badgepurple}{Por que $(-)\times(-)=(+)$?}{Intuição pelo padrão da tabuada}

\regra{Observe a sequência de $n \times (-2)$ à medida que $n$
decresce — cada passo \emph{sobe} $+2$:}

\exemp{%
  $3 \times (-2) = -6$\\
  $2 \times (-2) = -4$\\
  $1 \times (-2) = -2$\\
  $0 \times (-2) = \phantom{+}0$\\
  $-1 \times (-2) = +2$ \;{\footnotesize\color{mutedtext}(sobe $+2$)}\\
  $-2 \times (-2) = \hlm{+4}$ \;{\footnotesize\color{mutedtext}(neg.$\times$neg.$=$pos.)}%
}

\nota{\textbf{Interpretação:} multiplicar por $-1$ inverte o sinal.
Inverter duas vezes volta ao original:
$(-1)\times(-1)\times a = +a$.}

%─────────────────────────────────────────────────────────────
\TW{$+/+$}{darkblue}{Sinais Iguais → Positivo}{Pos.$\times$Pos. ou Neg.$\times$Neg.}

\formula{$(+)\times(+)=(+) \qquad (-)\times(-)=(+)$\\[3pt]
$(+)\div(+)=(+) \qquad (-)\div(-)=(+)$}

\exemp{%
  \textbf{Positivo $\times$ Positivo}\\
  $(+7)\times(+3) = +21$ \qquad $(+15)\div(+3) = \hlm{+5}$\\[3pt]
  \textbf{Negativo $\times$ Negativo}\\
  $(-7)\times(-3) = +21$ \qquad $(-15)\div(-3) = \hlm{+5}$%
}

\nota{\textbf{Regra prática:} sinais \textbf{iguais} → resultado
\emph{positivo}.}

%─────────────────────────────────────────────────────────────
\TW{$+/-$}{darkblue}{Sinais Diferentes → Negativo}{Pos.$\times$Neg. ou Neg.$\times$Pos.}

\formula{$(+)\times(-)=(-) \qquad (-)\times(+)=(-)$\\[3pt]
$(+)\div(-)=(-) \qquad (-)\div(+)=(-)$}

\exemp{%
  \textbf{Positivo $\times$ Negativo}\\
  $(+7)\times(-3) = -21$ \qquad $(+15)\div(-3) = \hlm{-5}$\\[3pt]
  \textbf{Negativo $\times$ Positivo}\\
  $(-7)\times(+3) = -21$ \qquad $(-15)\div(+3) = \hlm{-5}$%
}

\nota{\textbf{Regra prática:} sinais \textbf{diferentes} → resultado
\emph{negativo}.}

%─────────────────────────────────────────────────────────────
\TW{Alg}{darkblue}{Algoritmo de Cálculo}{Sinal primeiro, módulo depois}

\regra{%
\textbf{1.} Determine o sinal pela regra dos sinais.\\[2pt]
\textbf{2.} Multiplique (ou divida) os valores absolutos.\\[2pt]
\textbf{3.} Aplique o sinal ao resultado.}

\exemp{%
  \textbf{$(-12)\times(+5)$}\\
  P1: sinais diferentes → \textbf{negativo}\\
  P2: $12\times5 = 60$\\
  P3: $\hlm{-60}$ \ok%
}

\exemp{%
  \textbf{$(-48)\div(-6)$}\\
  P1: sinais iguais → \textbf{positivo}\\
  P2: $48\div6 = 8$\\
  P3: $\hlm{+8}$ \ok%
}

%─────────────────────────────────────────────────────────────
\TW{$n$ fat.}{badgegreen}{Múltiplos Fatores}{Conte os fatores negativos}

\regra{Em um produto com vários fatores, conte os negativos:\\[3pt]
\textbf{Quantidade par} de negativos → resultado \textbf{positivo}\\[2pt]
\textbf{Quantidade ímpar} de negativos → resultado \textbf{negativo}}

\exemp{%
  $(-2)\times(-3)\times(-4)$ — \textbf{3 neg.} (ímpar)\\
  $2\times3\times4 = 24 \;\Rightarrow\; \hlm{-24}$ \ok\\[3pt]
  $(-1)\times(-2)\times(-3)\times(-4)$ — \textbf{4 neg.} (par)\\
  $1\times2\times3\times4 = 24 \;\Rightarrow\; \hlm{+24}$ \ok\\[3pt]
  $(+5)\times(-2)\times(+3)\times(-1)$ — \textbf{2 neg.} (par)\\
  $5\times2\times3\times1 = 30 \;\Rightarrow\; \hlm{+30}$ \ok%
}

%─────────────────────────────────────────────────────────────
\TW{$a^n$}{darkblue}{Potências de Inteiros Negativos}{Par → positivo · Ímpar → negativo}

\formula{$(-a)^n = \begin{cases}
  +a^n & \text{se } n \text{ é par}\\
  -a^n & \text{se } n \text{ é ímpar}
\end{cases}$}

\exemp{%
  \textbf{Expoente par → positivo}\\
  $(-3)^2 = (-3)\times(-3) = +9$\\
  $(-2)^4 = +16$ \qquad $(-5)^2 = \hlm{+25}$\\[3pt]
  \textbf{Expoente ímpar → negativo}\\
  $(-3)^3 = -27$ \qquad $(-2)^5 = -32$\\
  $(-1)^{101} = \hlm{-1}$%
}

\nota{\textbf{O parêntese faz diferença:}\\
$(-2)^2 = +4$ (eleva $-2$ ao quadrado)\\
$-2^2 = -4$ (eleva $2$ e aplica o $-$ depois)}

%─────────────────────────────────────────────────────────────
\TW{Prop}{badgegreen}{Propriedades da Multiplicação}{Leis válidas para todos os inteiros}

\regra{%
\textbf{Comutativa:} $a\times b = b\times a$\\[2pt]
\textbf{Associativa:} $(a\times b)\times c = a\times(b\times c)$\\[2pt]
\textbf{Distributiva:} $a\times(b+c) = a\times b + a\times c$\\[2pt]
\textbf{Neutro:} $a\times 1 = a$\\[2pt]
\textbf{Absorvente:} $a\times 0 = 0$\\[2pt]
\textbf{Inversor de sinal:} $(-1)\times a = -a$}

\exemp{%
  \textbf{Distributiva:}\\
  $(-3)\times(4+2) = (-3)\times4 + (-3)\times2$\\
  $= -12 + (-6) = -18$\\
  Verificação: $(-3)\times6 = \hlm{-18}$ \ok%
}

%─────────────────────────────────────────────────────────────
\TW{$\checkmark$}{darkblue}{Relação $\times$ e $\div$ — Prova Real}{Divisão é a inversa da multiplicação}

\formula{Se $a\times b = c$, então $c\div b = a$ e $c\div a = b$}

\exemp{%
  $(-4)\times(+5) = -20$\\
  $(-20)\div(+5) = -4$ \ok\\
  $(-20)\div(-4) = \hlm{+5}$ \ok\\[3pt]
  $(-6)\times(-7) = +42$\\
  $(+42)\div(-7) = -6$ \ok\\
  $(+42)\div(-6) = \hlm{-7}$ \ok%
}

%─────────────────────────────────────────────────────────────
\TW{Tab}{badgegreen}{Tabela de Sinais}{Resumo das quatro combinações}

\begin{tcolorbox}[arc=5pt, boxrule=0.8pt, colback=white, colframe=vegreen,
  left=8pt, right=8pt, top=5pt, bottom=5pt]
{\color{vegreen}\bfseries\small SINAIS IGUAIS}%
\hfill{\footnotesize\color{vegreen}resultado positivo}\par\vspace{2pt}
{\small\color{bodytext}$(+)\times(+)=+$ \qquad $(-)\times(-)=+$}
\end{tcolorbox}\vspace{3pt}

\begin{tcolorbox}[arc=5pt, boxrule=0.8pt, colback=white, colframe=vered,
  left=8pt, right=8pt, top=5pt, bottom=5pt]
{\color{vered}\bfseries\small SINAIS DIFERENTES}%
\hfill{\footnotesize\color{vered}resultado negativo}\par\vspace{2pt}
{\small\color{bodytext}$(+)\times(-)=-$ \qquad $(-)\times(+)=-$}
\end{tcolorbox}

\obs{A mesma regra vale para a divisão: basta trocar $\times$ por
$\div$ em qualquer das linhas acima.}

%─────────────────────────────────────────────────────────────
\TW{Ex}{badgepurple}{Problemas Contextualizados}{Temperatura, finanças e movimento}

\exemp{%
  \textbf{P1 — Temperatura:} cai $3\,$°C por hora; em 5 horas?\\
  $(-3)\times(+5) = \hlm{-15}$ → cai $15\,$°C. \ok\\[3pt]
  \textbf{P2 — Dívida:} 4 dívidas de R\$\,12 cada.\\
  $4\times(-12) = \hlm{-48}$ → deve R\$\,48. \ok%
}

\columnbreak

\exemp{%
  \textbf{P3 — Compartilhar dívida:} R\$\,60 entre 5 pessoas.\\
  $(-60)\div(+5) = \hlm{-12}$ → cada um deve R\$\,12. \ok\\[3pt]
  \textbf{P4 — Movimento:} submarino a $-8\,$m/min; chega a $-56\,$m?\\
  $-56\div(-8) = \hlm{+7}$ → em $7$ minutos. \ok%
}

%─────────────────────────────────────────────────────────────
\TW{!}{vered}{Erros Comuns}{O que evitar nas provas}

\exemp{%
  \textbf{\nok\ Erro 1 — $(-)\times(-)=(-)$}\\
  $(-3)\times(-4) \neq -12$\\
  \ok\ Correto: sinais iguais → $\hlm{+12}$\\[4pt]
  \textbf{\nok\ Erro 2 — confundir $(-2)^2$ e $-2^2$}\\
  $(-2)^2 = +4$ \quad $-2^2 = -4$\\
  \ok\ O parêntese muda o resultado!%
}

\exemp{%
  \textbf{\nok\ Erro 3 — esquecer o sinal no quociente}\\
  $(-20)\div(+4) \neq 5$\\
  \ok\ Correto: sinais diferentes → $\hlm{-5}$\\[4pt]
  \textbf{\nok\ Erro 4 — contagem de negativos}\\
  $(-1)\times(-1)\times(-1) \neq +1$\\
  \ok\ 3 negativos (ímpar) → $\hlm{-1}$%
}

\end{multicols}

\end{document}
